% Part: proof-theory
% Chapter: natural-deduction
% Section: introduction

\documentclass[../../../include/open-logic-section]{subfiles}

\begin{document}

\olfileid[hi]{pt}{nat}{int}

\olsection{परिचय}

प्राकृतिक निगमन, !!{proof} प्रणालियों का ऐसा परिवार है जिसमें प्रत्येक
तार्किक संयोजक या परिमाणक के लिए अनुमानों का एक युग्म होता है: एक संकारक को
``प्रस्तुत'' करता है (!!a{introduction} नियम) और दूसरा संकारक का ``विलोपन''
करता है (!!a{elimination} नियम)। उदाहरणतः \Intro{\lif} के निष्कर्ष में~$!A
\lif !B$ रूप के !!{formula}s होते हैं और \Elim{\lif} के आधार-वाक्यों में
~$!A \lif !B$ रूप के !!{formula}s होते हैं।

प्राकृतिक निगमन के पहले दो रूप 1930 के दशक में गेरहार्ड गेंट्सेन और स्तानिस्वाव
याश्कोव्स्की ने प्रस्तुत किए। प्रमाण सिद्धांतकार अधिकतर गेंट्सेन की प्रणाली
की चर्चा करते हैं, जबकि याश्कोव्स्की से अध्यापन में सर्वाधिक प्रयुक्त
प्रणालियाँ उत्पन्न हुईं। दोनों प्रणालियों में कोई \emph{मान्यताओं} से आरंभ
कर सकता है---ऐसे !!{formula}s जिन्हें !!{prove} करने की आवश्यकता नहीं, पर
अन्य !!{formula}s को !!{proving} करने में प्रयोग किया जा सकता है। पूरा
!!{proof} ऐसी मान्यताओं पर निर्भर हो सकता है, और चूँकि वे अप्रमाणित हैं,
!!{proof} अपना निष्कर्ष (अपना \emph{अंतिम-!!{formula}}) स्थापित नहीं करता,
बल्कि केवल यह स्थापित करता है कि उसका निष्कर्ष मान्यताओं से निष्पन्न होता है।

कुछ अनुमान नियम ऐसे हैं जिनके निष्कर्ष अब कुछ मान्यताओं पर निर्भर नहीं रहते:
वे मान्यताओं को ``मुक्त'' करते हैं। उदाहरणतः \Intro{\lif} नियम, मान्यता~$!A$
से निष्कर्ष~$!B$ का !!a{proof} लेता है और उसका निष्कर्ष $!A \lif !B$ होता
है। $!A \lif !B$ पर समाप्त होने वाला !!{proof} अब~$!A$ पर निर्भर नहीं है;
मान्यता $!A$ मुक्त कर दी जाती है। गेंट्सेन की प्रणालियों में !!{proof},
!!{formula}s के वृक्ष होते हैं, मान्यताएँ सबसे ऊपर के !!{formula}s होती हैं,
और~\Intro{\lif} जैसे नियमों पर ऐसे चिह्न होते हैं जो बताते हैं कि कौन-सी
मान्यताएँ मुक्त की गई हैं। याश्कोव्स्की की प्रणालियों में किसी मान्यता पर
निर्भर !!{proof} के अंशों को किसी ढंग से अलग किया जाता है, जैसे खड़ी रेखा के
पीछे भीतर करके या डिब्बे में; मुक्त मान्यता~$!A$ और सशर्त का अनुवर्ती $!B$
डिब्बे की पहली और अंतिम पंक्तियाँ होते हैं, तथा \Intro{\lif} का निष्कर्ष~$!A
\lif !B$ डिब्बे के बाहर होता है।

ये प्रणालियाँ ``प्राकृतिक'' हैं क्योंकि अनुमान नियम उस ढंग को प्रतिबिंबित
करते हैं जिससे हम (या कम-से-कम गणितज्ञ) स्वाभाविक रूप से तर्क करते हैं। किसी
काल्पनिक परिणाम (सशर्त) को स्थापित करने के लिए हम पूर्ववर्ती मानकर उससे
अनुवर्ती सिद्ध करते हैं। यही~\Intro{\lif} है। जब हमें कोई सशर्त $!A \lif !B$
और उसका पूर्ववर्ती~$!A$ ज्ञात हो, तो हम~$!B$ निष्कर्षित करते हैं। यही
~\Elim{\lif} है। यह दिखाने के लिए कि किसी वियोजन~$!A \lor !B$ को मानने पर
कोई निष्कर्ष सत्य है, हम स्थितियों द्वारा प्रमाण प्रयोग करते हैं। यही
~\Elim{\lor} है। सामान्य दावा~$\lforall[x][!A(x)]$ सिद्ध करने के लिए हम किसी
मनमाने~$c$ के लिए $!A(c)$ सिद्ध करते हैं। यही~\Intro{\lforall} है। और इसी
प्रकार आगे।

उदाहरणतः गेंट्सेन-शैली के प्राकृतिक निगमन में~$!A \lif (!A \lor !B)$ का यह
एक सरल !!{proof} है:
\[
\AxiomC{$\Discharge{!A}{x}$}
\RightLabel{\Intro{\lor}}
\UnaryInfC{$!A \lor !B$}
\DischargeRule{\Intro{\lif}}{x}
\UnaryInfC{$!A \lif (!A \lor !B)$}
\DisplayProof
\]
यह मान्यता~$!A$ से आरंभ होता है,~$!A \lor !B$ अनुमानित करने के लिए
\Intro{\lor}, और फिर~$!A \lif (!A \lor !B)$ अनुमानित करने के लिए
\Intro{\lif} का उपयोग करता है। \Intro{\lif} अनुमान मान्यता~$!A$ को मुक्त
करता है; इसे चिह्न~$x$ बताता है।

प्रमाण सिद्धांतकार !!{formula}s के वृक्षों पर काम करने वाली गेंट्सेन की मूल
प्रणालियाँ पसंद करते हैं, यद्यपि एक रूपांतर भी प्रमाण सिद्धांत में महत्त्वपूर्ण
है। मान्यताओं~$\Gamma$ से~$!A$ का !!^a{proof} दिखाता है कि $!A$,~$\Gamma$
से निष्पन्न होता है, अर्थात् $\Gamma \Entails !A$। प्राकृतिक निगमन के अनुमान
नियमों को स्वाभाविक रूप से ऐसे तात्पर्यों के बारे में कुछ बताते हुए पढ़ा जा
सकता है, उदाहरणतः \Intro{\lif} कहता है कि यदि $\Gamma + !A \Entails !B$,
तो $\Gamma \Entails !B$। नियमों को सीधे इस तरह सूत्रबद्ध करना स्वाभाविक है
कि वे अनुमान नियम के आधार-वाक्यों और निष्कर्ष, दोनों में मान्यताओं तथा उनसे
निष्पन्न निष्कर्ष पर काम करें, अर्थात् वे अनुक्रमों $\Gamma \Sequent !A$ पर
काम करें। ऐसी प्रणाली में ऊपर का सरल प्रमाण इस तरह दिखेगा:
\[
\Axiom$x:!A \fCenter !A$
\RightLabel{\Intro{\lor}}
\UnaryInf$x:!A \fCenter !A \lor !B$
\DischargeRule{\Intro{\lif}}{x}
\UnaryInf$\fCenter !A \lif (!A \lor !B)$
\DisplayProof
\]
जैसा आप देखते हैं, नियम वही हैं: वे~$\Sequent$ के दाईं ओर के !!{formula} पर
काम करते हैं; बाईं ओर के !!{formula}s केवल उन मान्यताओं को सूचीबद्ध करते हैं
जिन पर वे प्रत्येक चरण में निर्भर हैं।

!!{introduction}/!!{elimination} नियम-युग्मों का पूरा समुच्चय अपने आप में
शास्त्रीय तर्कशास्त्र के लिए पूर्ण नहीं है। हमें~$\lfalse$ से संबंधित दो
अतिरिक्त नियम जोड़ने होंगे। पहला ``अंतर्ज्ञानवादी असंगति नियम''~\FalseInt,
या ``विस्फोट'', है जो कहता है कि $\lfalse$ से हम कुछ भी अनुमानित कर सकते हैं।
दूसरा अप्रत्यक्ष प्रमाण का शास्त्रीय सिद्धांत~\FalseCl{} (या ``शास्त्रीय
असंगति नियम'') है, जो कहता है कि यदि हम~$\lnot !A$ से~$\lfalse$ को
!!{prove} कर सकते हैं, तो हमने~$!A$ को !!{prove} कर दिया। यदि~\FalseCl को
छोड़ दें, तो अंतर्ज्ञानवादी तर्कशास्त्र के लिए उपयुक्त प्रणाली मिलती है। यदि
दोनों को छोड़ दें, तो वह ``न्यूनतम तर्कशास्त्र'' कहलाता है।

हम शास्त्रीय और अंतर्ज्ञानवादी तर्कशास्त्र के लिए गेंट्सेन-शैली के प्राकृतिक
निगमन की चर्चा करेंगे। ये प्रणालियाँ \Log{N1c} और \Log{N1i} हैं। प्रणालियाँ
\Log{N2c} और \Log{N2i} इनकी संगत प्रणालियाँ हैं, जो अकेले !!{formula}s~$!A$
के स्थान पर अनुक्रमों $\Gamma \Sequent !A$ पर काम करती हैं।

\end{document}
